A bőrrák, és különösen a melanoma, a korai felismerés esetén jól kezelhető, késői
stádiumban azonban drámaian romlik a túlélési esély; a bőrgyógyászati
szakellátáshoz való hozzáférés ugyanakkor világszerte egyenetlen. Jelen dolgozat
egy végponttól végpontig megvalósított, mesterséges intelligencián alapuló
bőr\-elváltozás-elemző rendszert (DermAI) mutat be, amely a korai szűrést egy
okostelefonon és webböngészőben egyaránt elérhető eszközzel támogatja.

A dolgozat kutatási része egy Vision Transformer (ViT) modellt tanít be az ISIC
2024 adathalmazon a bőrelváltozások jó\-indulatú és rosszindulatú kategóriába
sorolására, kiemelt figyelmet fordítva az adathalmaz erős
osztály-egyensúlytalanságának kezelésére. A modellt egy azonos feltételek mellett
betanított ResNet-18 konvolúciós hálóval és egy zero-shot Gemini nagy nyelvi
modellel hasonlítjuk össze, a hivatalos ISIC 2024 metrikán, statisztikai
szignifikancia-vizsgálattal és kalibrációs elemzéssel kiegészítve. A mérnöki rész
a betanított modellt egy teljes alkalmazássá egészíti ki: mobil- és webklienssel,
a modellt kiszolgáló MCP-réteggel, a Gemini modellre épülő képvalidációval és
természetes nyelvű magyarázattal, valamint hitelesített, felhasználónként
elkülönített felhőalapú adattárolással.

\textbf{Kulcsszavak}: bőrrák, melanoma, mély tanulás, Vision Transformer, ISIC
2024, osztály-egyensúlytalanság, nagy nyelvi modell, mobil egészségügy
